\documentclass[11pt,a4paper]{article}
\usepackage[T1]{fontenc}
\usepackage[utf8]{inputenc}
\usepackage[main=italian,provide=*]{babel}
\usepackage{amsmath,amssymb}
\usepackage{geometry}
\geometry{margin=2.5cm}
\usepackage{booktabs}
\usepackage{seqsplit}
\usepackage{hyperref}
\hypersetup{colorlinks=true,linkcolor=blue,citecolor=blue,urlcolor=blue}

\title{Benchmark esatto (Fase 3) — trimero ad anello}
\author{Note di lavoro — Tesi triennale, Samuele Schiavi\\ Relatore: Prof.\ Alessandro Chiesa}
\date{}

\begin{document}
\maketitle

\section*{Introduzione}

Questo documento raccoglie i moduli di benchmark esatto per il trimero di spin-1/2
($N=3$) in topologia ad anello (triangolo isoscele). Sono i moduli su cui tutto il
resto del lavoro (VQE, e in futuro Trotter e correlazioni) si appoggia senza mai
modificarli: la diagonalizzazione numerica (\texttt{eigh}) è il metro di misura con
cui ogni risultato variazionale viene giudicato, esattamente come
\texttt{dimer\_exact.py} lo è per il dimero.

Per ciascun file: cosa fa, con quali verifiche è stato controllato in questa sessione,
e cosa se ne può concludere — per i moduli \texttt{.py}, anche le funzioni interne e
come si usa il file.

\textbf{Nota.} \texttt{verifica.py} (ora rinominato \texttt{verifica\_anello.py}) è
stato spostato nel documento dedicato \texttt{trimero\_anello\_verifiche.pdf}, insieme
al suo gemello \texttt{verifica\_stati\_anello.py} — non è più descritto qui.

\section*{\texttt{trimer\_ring\_exact.py}}

Modulo di benchmark: Hamiltoniana come \texttt{SparsePauliOp}
($H=J\,\boldsymbol\sigma_1\cdot\boldsymbol\sigma_2+J'(\boldsymbol\sigma_2\cdot\boldsymbol\sigma_3+\boldsymbol\sigma_3\cdot\boldsymbol\sigma_1)+b\sum_iZ_i$),
spettro analitico via decomposizione di Kambe (simmetria di scambio
$1\leftrightarrow2$, conservazione di $\mathbf S_{12}^2$), gli otto stati di Kambe
espliciti in base computazionale, il campo critico $b_c(J,J')$ nei diversi regimi di
segno, e le funzioni per il termine DM (\texttt{trimer\_hamiltonian\_dm},
\texttt{exact\_sweep\_dm}, \texttt{ground\_state\_projector\_dm} — disponibili nel
modulo esatto ma non usate nella fase VQE senza DM).

\textbf{Verifiche.} Otto self-test a precisione macchina (spettro vs formula analitica
su punti casuali, stati di Kambe come autostati esatti verificati contro l'operatore
Qiskit, ortonormalità, numeri quantici, simmetria $P_{12}$, proiettore sul
fondamentale, campo critico teorico vs numerico, magnetizzazione mediata sul
sottospazio degenere) più tre self-test dedicati al DM. Rieseguiti in questa sessione:
tutti superati, incluso dopo la correzione applicata a $\langle M_z\rangle(b{=}0)$
(media sul sottospazio degenere, non un singolo autovettore arbitrario).

\textbf{Considerazioni.} Il file ha subito una correzione reale durante il lavoro:
prima della sessione corrente, $\langle M_z\rangle$ a $b=0$ (fondamentale degenere)
veniva calcolato su un singolo autovettore restituito da \texttt{numpy.linalg.eigh},
scelto arbitrariamente fra quelli del sottospazio degenere — il salto apparente da
$+0.5$ a $-0.5$ osservato in una vecchia figura era un artefatto numerico, non un
effetto fisico. Il salto \emph{vero} (fisico, discontinuità del ground state) è solo
al campo critico $b_c$.

\textbf{Conclusione.} Modulo maturo, verificato più volte in sessioni diverse con
esito sempre positivo. È la base condivisa da ogni file VQE dell'anello.

\textbf{Funzioni interne.}
\begin{itemize}
\item \texttt{trimer\_hamiltonian(J,\ Jp,\ b)} — l'Hamiltoniana completa come
\texttt{SparsePauliOp}.
\item \texttt{magnetization\_operator()} — $M_z=(Z_1+Z_2+Z_3)/2$, convenzione di
spin ($s=\sigma/2$).
\item \texttt{S12\_squared\_operator()} — il Casimir della coppia base,
$S_{12}^2=\tfrac32 I+\tfrac12\boldsymbol\sigma_1\cdot\boldsymbol\sigma_2$.
\item \texttt{S\_total\_squared\_operator()} — il Casimir totale $S^2$.
\item \texttt{analytic\_eigenvalues(J,\ Jp,\ b)} — le otto energie esatte
$E(S_{12},S,M)=E_\text{blocco}+2bM$, ordinate.
\item \texttt{critical\_field(J,\ Jp)} — il campo critico $b_c$ dell'incrocio del
fondamentale, a seconda di chi è il ground state a $b=0$.
\item \texttt{kambe\_states()} — dizionario \texttt{\{(blocco, M): vettore a 8
componenti\}} con gli otto stati espliciti.
\item \texttt{kambe\_energy(J,\ Jp,\ b,\ block,\ M)} — l'energia analitica del
singolo stato \texttt{(block, M)}.
\item \texttt{exact\_sweep(b\_values,\ J,\ Jp)} — diagonalizza $H$ su una griglia di
campo $b$, restituisce spettro completo e osservabili del fondamentale.
\item \texttt{ground\_state\_projector(J,\ Jp,\ b,\ tol)} — il proiettore sul
multipletto fondamentale, robusto a eventuale degenerazione.
\item \texttt{dm\_term(i,\ j,\ D)} — il termine DM $D(X_iZ_j-Z_iX_j)$ su una coppia
di siti (stessa forma già validata per il dimero).
\item \texttt{trimer\_hamiltonian\_dm(J,\ Jp,\ b,\ mode,\ D)} — l'Hamiltoniana con
il termine DM aggiunto, secondo l'opzione scelta (\texttt{mode}).
\item \texttt{dm\_min\_gap(J,\ Jp,\ mode,\ D,\ \seqsplit{search\_half\_width})} — il
gap vero fra i due autovalori più bassi, minimizzato su tutto $b$ (non calcolato a un
valore di campo fissato — la stessa distinzione metodologica cruciale di
\texttt{analisi\_dm\_trimero.pdf}).
\item \texttt{exact\_sweep\_dm(b\_values,\ J,\ Jp,\ mode,\ D)} — come
\texttt{exact\_sweep}, ma con l'Hamiltoniana DM.
\item \texttt{ground\_state\_projector\_dm(J,\ Jp,\ b,\ mode,\ D,\ tol)} — come
\texttt{ground\_state\_projector}, ma con l'Hamiltoniana DM.
\item \texttt{\_self\_test()} — gli otto self-test base, eseguiti automaticamente se
il file viene lanciato come script.
\item \texttt{\_self\_test\_dm()} — i tre self-test dedicati al DM.
\end{itemize}

\textbf{Come si usa.} Come libreria: \texttt{from trimer\_ring\_exact import
trimer\_hamiltonian, exact\_sweep, \ldots} — è la dipendenza comune a tutti i moduli
VQE dell'anello. Come script standalone: \texttt{python3 trimer\_ring\_exact.py}
esegue entrambi i blocchi di self-test e rigenera la figura
\texttt{trimer\_ring\_exact.png}.

\section*{\texttt{trimer\_ring\_exact.png}}

Figura generata dal modulo precedente: spettro completo e magnetizzazione del
fondamentale, con annotazione esplicita che distingue il punto convenzionale $b=0$
(fondamentale degenere, valore di $\langle M_z\rangle$ dipendente dalla convenzione di
media) dall'incrocio fisico vero a $b_c$ (discontinuità reale). Rigenerata dopo la
correzione sopra descritta.

\section*{\texttt{teoria\_trimero\_isoscele.pdf}}

Documento teorico: dimostrazione (non solo enunciato) delle identità
$\sigma_1\cdot\sigma_2=2P_{12}-I$ e $S_{12}^2=I+P_{12}$, prova esplicita che
$[H,S_{12}^2]=0 \Leftrightarrow J'_{23}=J'_{31}$, decomposizione di Kambe con conteggio
dei multipletti, spettro chiuso, mappa dei segni completa nel piano $(J,J')$, campo
critico nei diversi regimi, dimostrazione rigorosa (non solo numerica) che gli
incroci sono veri, non anticrossing.

\textbf{Conclusione.} È la sorgente teorica primaria da cui discende tutto il resto:
\texttt{trimer\_ring\_exact.py} implementa esattamente quanto derivato qui.

\section*{\texttt{derivazione\_stati\_kambe\_trimero.pdf}}

Derivazione esplicita degli otto autostati nella base computazionale
$|q_1q_2q_3\rangle$, via operatore di abbassamento $S_-$ per il blocco $A$ e
ortogonalità per il blocco $B$. È il documento da cui \texttt{kambe\_states()} in
\texttt{trimer\_ring\_exact.py} è stato trascritto in codice, e da cui deriva
\seqsplit{\texttt{\_prepare\_initial\_state}} nei moduli VQE (i pesi
$1/\sqrt6,1/\sqrt6,-\sqrt{2/3}$ del blocco $B$ vengono direttamente da qui).

\textbf{Conclusione.} Stesso metodo poi riusato, con gli adattamenti necessari, per
derivare gli otto stati della catena aperta (documento gemello per la catena).

\section*{\texttt{animazione\_trimero\_isoscele.html}}

Visualizzazione interattiva: geometria del triangolo, mappa di fase 2D nel piano
$(J,J')$, spettro con il fondamentale colorato per blocco. Utile per l'intuizione
qualitativa (dove si sposta il fondamentale cambiando $J,J'$), non per i numeri
precisi (quelli vengono dal modulo esatto).

\section*{Conclusione generale}

Il modulo esatto dell'anello ha lo stesso livello di rigore già stabilito per il
dimero: self-test a precisione macchina, verifica indipendente su griglia estesa,
correzione tracciata quando trovata (media sul sottospazio degenere). È pronto a
supportare tutto il lavoro VQE già costruito e le fasi ancora da affrontare (Trotter,
correlazioni dinamiche).

\end{document}
